% s-210-p210-m003; original PDF 210, 211
\Aphorism{s-210-p210-m003}{304}
Das ideale Drama zu n Personen muss enthalten: n Persönlichkeiten, \Math{\frac{\mathrm{n}(\mathrm{n}-1)}{2}} dialogische Beziehungen (Unterdramen zu je Zweien), \Math{\frac{\mathrm{n}(\mathrm{n}-1)(\mathrm{n}-2)}{6}} „dreieckige“ Verhältnisse (Unterdramen zu je Dreien) und so fort, oder, wie man in der Arithmetik sagt, Unionen, Binionen, Ternionen in der überhaupt möglichen Anzahl. Also innere Vollständigkeit und Erschöpfung aller denkbaren Combinationen; keiner der Handelnden soll nur einseitig wirken oder einseitig empfangen — zwischen den vorhandenen Punkten sind alle Linien, Ebenen u. s. w. wirklich zu zeichnen. Daneben muss das Drama ein Gesamtvorgang zwischen allen n Personen, ja womöglich Theil eines Vorgangs zwischen mehr als n Personen sein, d. h. es muss einen Hintergrund (ein Milieu) geschichtlicher oder sonst allgemeiner Art haben. Man sieht, wie sehr der ernsthafte Dramatiker bemüht sein wird, sein n zu \Sperr{verkleinern}!
\par

